\documentclass[main.tex]{subfiles}


\begin{document}

% \AddToShipoutPictureFG{%
% \begin{tikzpicture}[overlay,remember picture]
% \path (current page.south west) -- (current page.north east)
% node[midway,scale=1,color=gray,sloped,opacity=0.4] {\textnormal{Кравченко Игорь Игоревич\hspace{1cm} +79010144910\hspace{1cm} super.antig@yandex.ru}};
% \end{tikzpicture}
% }

% %from pagenumber to up
% \phantomsection 
% \hypertarget{toc}{}

 

 

 
   

\section{Сложение и относительность скоростей}


Движение тела, разумеется, можно рассматривать не только относительно планеты. На рис.\,\ref{fig:p13} показан опыт с катящимся шаром по неподвижной тележке.

\begin{figure}[H]
    \centering

    \begin{tikzpicture}[font=\small,            line cap=round,                             >={Stealth[inset=0pt,round,length=3mm, angle'=20]},vec/.style={>={Stealth[inset=0pt,round,length=3.5mm, angle'=20]}}]


\filldraw[lightgray,semitransparent] (0,0) rectangle (12,-0.3);
\draw (0,0) -- (12,0);

\begin{scope}[xshift=4cm]    
%%motion1
\filldraw[lightgray,draw=black] (0,0.5) rectangle (4,1);
\node [left,black] at (0,.75) {Т};
% PERSON
\def\W{5.2}     % ground width
\def\L{3.2}     % length
\def\T{0}    % plank thickness
\def\H{2.0}     % human height
  \def\x{1} % human position
  \coordinate (O) at (0,0);
\begin{scope}[yshift=1cm]  
  \draw[thick] (\x,\H) circle(0.3) coordinate (H) node[xshift=0.3cm,right] {Н};
  \draw[thick] (H)++(-90:0.3) coordinate (N) to[out=-92,in=92]++ (0,-0.40*\H) coordinate (P);
  \draw[thick] (N)++(-95:0.03) to[out=-105,in=90]++ (-0.02*\W,-0.4*\H);
  \draw[thick] (N)++(-85:0.03) to[out=-80,in=90]++ (0.02*\W,-0.4*\H);
  \draw[thick] (P) to[out=-92,in=82] (\x-0.02*\W,\T);
  \draw[thick] (P) to[out=-80,in=90] (\x+0.02*\W,\T);
\end{scope}

%wheels
\draw[black,fill=black]
    (1,.3) circle(.3cm)
    (3,.3) circle(.3cm);
\draw[fill=gray]
    (1,.3) circle(.2cm)
    (3,.3) circle(.2cm);

%%%%%%%%%%%%%%%%%%%%%%%%%%%%%%%%%%
%ball
\shade[ball color=lightgray,nearly transparent] (2,1.3) circle (.3);
\shade[ball color=lightgray,semitransparent,xshift=1cm] (2,1.3) circle (.3);
\shade[ball color=lightgray,xshift=2cm] (2,1.3) circle (.3) node[above,yshift=.3cm,black] {Ш} node [circle,fill=NavyBlue,inner sep=0pt,minimum size=1mm,opaque] (vo) {};
\draw[->,NavyBlue,thick,vec] (vo) --++ (0:1) node[black,above right,inner xsep=0] {$\vec v_\text{отн}$};
\end{scope}



    \end{tikzpicture}
\caption{Тележка покоится}
\label{fig:p13}
\end{figure}

% \begin{figure}[H]
%     \centering

%     \includestandalone{PICS/mech/7add_1}

% \caption{Тележка покоится}
% \label{fig:p13}
% \end{figure}


Наблюдатель~Н толкает шар~Ш вправо и измеряет скорость шара~$\vec v_\text{отн}$ относительно тележки~Т. Такой толчок дал бы ту же скорость, даже если бы тележка катилась равномерно прямолинейно: ведь наблюдатель стоит на ней.

Ясно, что если шар толкают так же, а тележка движется, то расстояние пройденное шаром за то же время относительно планеты будет другим (рис.\,\ref{fig:p14}). 

%%%%%%%%%%%%%%%%%%%%%%%%%%%%%%%%%%%%
% the next pic
%%%%%%%%%%%%%%%%%%%%%%%%%%%%%%%%%%%%

\begin{figure}[H]
    \centering

    \begin{tikzpicture}[font=\small,            line cap=round,                             >={Stealth[inset=0pt,round,length=3mm, angle'=20]},vec/.style={>={Stealth[inset=0pt,round,length=3.5mm, angle'=20]}}]


\filldraw[lightgray,semitransparent] (0,0) rectangle (12,-0.3);
\draw (0,0) -- (12,0);

% PERSON
\def\W{5.2}     % ground width
\def\L{3.2}     % length
\def\T{0}    % plank thickness
\def\H{2.0}     % human height
  \def\x{1} % human position
  \coordinate (O) at (0,0);
\begin{scope}[yshift=0cm]  
  \draw[thick] (\x,\H) circle(0.3) coordinate (H) node[xshift=0.3cm,right] {Н};
  \draw[thick] (H)++(-90:0.3) coordinate (N) to[out=-92,in=92]++ (0,-0.40*\H) coordinate (P);
  \draw[thick] (N)++(-95:0.03) to[out=-105,in=90]++ (-0.02*\W,-0.4*\H);
  \draw[thick] (N)++(-85:0.03) to[out=-80,in=90]++ (0.02*\W,-0.4*\H);
  \draw[thick] (P) to[out=-92,in=82] (\x-0.02*\W,\T);
  \draw[thick] (P) to[out=-80,in=90] (\x+0.02*\W,\T);
\end{scope}

%all mot
\begin{scope}[xshift=4cm]
%scope for motions 1 and 2
\begin{scope}[transparency group, nearly transparent]
    
%%%%%%%%%%%%%%%%%%%%%%%%%%%%%%%%%%%%%%%
\begin{scope}[transparency group, semitransparent]
%%motion1
\filldraw[lightgray,draw=black] (0,0.5) rectangle (4,1);


%wheels
\draw[black,fill=black]
    (1,.3) circle(.3cm)
    (3,.3) circle(.3cm);
\draw[fill=gray]
    (1,.3) circle(.2cm)
    (3,.3) circle(.2cm);
\end{scope}

%%%%%%%%%%%%%%%%%%%%%%%%%%%%%%%%%%%%%%%
\begin{scope}[xshift=1cm,transparency group]
    
%%motion1
\filldraw[lightgray,draw=black] (0,0.5) rectangle (4,1);


%wheels
\draw[black,fill=black]
    (1,.3) circle(.3cm)
    (3,.3) circle(.3cm);
\draw[fill=gray]
    (1,.3) circle(.2cm)
    (3,.3) circle(.2cm);
\end{scope}
\end{scope}



%%%%%%%%%%%%%%%%%%%%%%%%%%%%%%%%%%%%%%%
\begin{scope}[xshift=2cm]
    
%%motion1
\filldraw[lightgray,draw=black] (0,0.5) rectangle (4,1);
\node [left,black] at (0,.75) {Т};
\draw[->,NavyBlue,thick,vec] (4,.75) node [circle,fill=NavyBlue,inner sep=0pt,minimum size=1mm,opaque] {} --++ (0:1) node[black,below right,inner xsep=0] {$\vec v_\text{пер}$};
;

%wheels
\draw[black,fill=black]
    (1,.3) circle(.3cm)
    (3,.3) circle(.3cm);
\draw[fill=gray]
    (1,.3) circle(.2cm)
    (3,.3) circle(.2cm);
\end{scope}
\end{scope}

%ball
\shade[ball color=lightgray,nearly transparent] (6,1.3) circle (.3);
\shade[ball color=lightgray,semitransparent] (8,1.3) circle (.3);
\shade[ball color=lightgray] (10,1.3) circle (.3) node[above,yshift=.3cm,black] {Ш} node [circle,fill=NavyBlue,inner sep=0pt,minimum size=1mm,opaque] (va) {};
\draw[->,NavyBlue,thick,vec] (va) --++ (0:2) node[black,above right,inner xsep=0] {$\vec v_\text{абс}$};
;


    \end{tikzpicture}
\caption{Тележка движется}
\label{fig:p14}
\end{figure}

% \begin{figure}[H]
%     \centering

%     \includestandalone{PICS/mech/7add_2}

% \caption{Тележка движется}
% \label{fig:p14}
% \end{figure}


Шар снова толкают вправо и его скорость в системе отсчета тележки опять равна~$\vec v_\text{отн}$ (на рис.\,\ref{fig:p14} не показана). В этот раз тележка имеет скорость~$\vec v_\text{пер}$ и как бы переносит тело. В системе отсчета планеты шар теперь двигается с некоторой скоростью~$\vec v_\text{абс}$, называемой \emph{абсолютной скоростью}. Cвязь трех ключевых скоростей, упомянутых ранее, дает \textbf{закон сложения скоростей}:
\begin{equation}
    \vec v_\text{абс}=\vec v_\text{отн}+\vec v_\text{пер},
\end{equation}
где $\vec v_\text{отн}$~--- \emph{относительная скорость}, $\vec v_\text{пер}$~--- \emph{переносная скорость}.



% \begin{law}
% \textbf{Закон сложения скоростей:}
% \begin{equation}
%     \vec v_\text{абс}=\vec v_\text{отн}+\vec v_\text{пер}
% \end{equation}
% где $\vec v_\text{отн}$~--- относительная скорость; $\vec v_\text{пер}$~--- переносная скорость.
% \end{law}

% Итак, \emph{относительность скоростей} заключается в том, что необходимо пояснять для какой системы отсчета указывается та или иная скорость.

Иногда решение задачи становится проще, если рассматривать движение одного тела, прикрепившись к другому. На рис.\,\ref{fig:p15} красное тело догоняет синее.

    \begin{figure}[H]
    \centering

    \begin{tikzpicture}[font=\small,            line cap=round,                             >={Stealth[inset=0pt,round,length=3mm, angle'=20]},vec/.style={>={Stealth[inset=0pt,round,length=3.5mm, angle'=20]}}]


\filldraw[lightgray,semitransparent] (0,0) rectangle (12,-0.3);
\draw (0,0) -- (12,0);

\shade[ball color=red] (1,.3) circle (.3);
\shade[ball color=NavyBlue] (6,.3) circle (.3);

%vec
\draw[->,NavyBlue,thick,vec] (1.4,.3) --++ (2,0) node[black,above] {$\vec v_1$};

\draw[->,NavyBlue,thick,vec] (6.4,.3) --++ (1,0) node[black,above] {$\vec v_2$};


    \end{tikzpicture}
\caption{Красный шар настигает синий шар}
\label{fig:p15}
\end{figure}

% \begin{figure}[H]
%     \centering

%     \includestandalone{PICS/mech/7add_3}

% \caption{Красный шар настигает синий шар}
% \label{fig:p15}
% \end{figure}


Если стоять на синем (правом) шаре, скользящем, например, без вращения, можно измерять скорость сближения~--- то есть \textbf{относительную скорость}: 
\begin{equation}
    \vec v_{12}=\vec v_{1}-\vec v_{2},
\end{equation}
где $\vec v_\text{12}$~--- скорость первого тела относительно второго тела.





 
\end{document}